\section{CAT3D: modelos de difusión multi-vida + destilación NeRF}

\subsection{Resumen del método}

CAT3D es un método que acepta como entrada texto, una sola imagen o múltiples imágenes para generar escenas 3D completas. Su idea central combina modelos de difusión multi-vidad con el flujo de reconstrucción NeRF.

\subsection{Arquitectura del modelo}

La entrada de CAT3D incluye dos tipos de vistas:

\begin{itemize}
    \item \textbf{Vistas observadas (Observed Views)}: incluyen la imagen de entrada y sus parámetros de cámara correspondientes. La imagen se codifica mediante un VAE preentrenado en variables latentes de imagen, mientras que la información de la cámara se codifica como un mapa de coordenadas de rayo por píxel, que contiene el origen y la dirección del rayo para cada píxel.
    \item \textbf{Vistas nuevas (Novel Views)}: no tienen imágenes reales; en su lugar, cuentan con variables latentes de ruido, acompañadas de información sobre la pose de la cámara para la vista objetivo.
\end{itemize}

Para ayudar al modelo a distinguir entre ambos tipos de vistas, también se introduce un \textbf{máscara}, que indica cuáles son las vistas observadas y cuáles están pendientes de generación.

\begin{knowledgebox}{Mapa de coordenadas de rayo (Ray Coordinate Map)}
Un mapa de coordenadas de rayo es una forma de codificar la información de la pose de la cámara como una representación espacial. Para cada píxel en la imagen, se calcula su origen de rayo y su dirección de rayo correspondientes; esta información se almacena como un mapa de coordenadas del mismo tamaño que la imagen. Esta representación permite que el modelo entienda directamente las relaciones geométricas 3D a nivel de píxel.
\end{knowledgebox}

El modelo utiliza \textbf{mecanismos de atención} entre todos los marcos, lo que le permite generar salidas consistentes en múltiples vistas. Los datos de entrenamiento consisten en grandes cantidades de imágenes multi-vidad con anotaciones de pose, supervisando al modelo para generar nuevas vistas consistentes dadas las vistas conocidas.

\subsection{Flujo de dos etapas}

CAT3D adopta un flujo de dos etapas:

\begin{enumerate}
    \item \textbf{Generación multi-vidad}: dada una pequeña cantidad de vistas de entrada, se utiliza un modelo de difusión multi-vidad para generar un conjunto \textbf{denso} de vistas de la escena.
    \item \textbf{Reconstrucción NeRF}: todas las vistas generadas se introducen en el pipeline estándar de NeRF para realizar una reconstrucción 3D.
\end{enumerate}

\begin{warningbox}{Las múltiples vistas generadas no son consistentes con 3D perfectas}
Las imágenes generadas por modelos de difusión multi-vidad no son completamente consistentes en 3D; pueden existir pequeñas inconsistencias entre diferentes vistas. Sin embargo, la clave está en que estas imágenes son \textbf{suficientemente consistentes} para que el pipeline de reconstrucción NeRF pueda integrarlas en una salida 3D coherente. En esencia, NeRF actúa como un papel de «promedio», eliminando las inconsistencias presentes en las múltiples vistas.
\end{warningbox}

\subsection{Rendimiento}

De manera end-to-end, CAT3D puede generar una escena 3D completa a partir de cualquier cantidad de entradas (incluso solo una imagen) en aproximadamente \textbf{1 minuto}.

\subsection{Resumen de la sección}

CAT3D valida la viabilidad del paradigma de «generación 2D + destilación 3D»: los modelos de difusión multi-vidad se encargan de generar imágenes de alta calidad, mientras que NeRF se encarga de destilar esas imágenes en representaciones 3D consistentes. Sin embargo, el paso de optimización de NeRF sigue siendo el cuello de botella para la velocidad de inferencia.

